\chapter*{Abstract}
\addcontentsline{toc}{chapter}{Abstract}
\markright{Abstract}

Fine-tuning one pre-trained language model per task gives strong task-specific performance, but it multiplies storage and deployment effort, while static weight merging discards much of the knowledge gained during fine-tuning. This report studies Weight-Ensembling Mixture of Experts (WEMoE) as a model consolidation method for natural language processing. Five BERT-base experts are fine-tuned independently on the MRPC, RTE, CoLA, SST-2 and QNLI tasks of the GLUE benchmark. In the merged model, non-MLP parameters are combined statically through Task Arithmetic, whereas the MLP parameters of every encoder layer are reconstructed for each input as the pre-trained weights plus a router-weighted sum of frozen task vectors. Predictions therefore come from a single forward pass with parameter-level combination and a task-specific head, not from averaged logits. Unlike the original image-based WEMoE study, which adapts routers on unlabelled test data, this implementation trains only the routers with Few-Shot Fine-Tuning on labelled examples that the experts never saw. All model selection uses a dedicated development split, and final scores are computed on the official validation splits. With 256 examples per task, WEMoE reaches an average score of $69.9\pm0.1$, against 57.4 for Task Arithmetic, 78.2 for the experts and 45.8 for a majority-class predictor. It recovers 59\% of the gain between the majority baseline and the experts. Upscaling only the first MLP layer gives a similar average ($69.1\pm0.1$) with 47.2\% of the storage of five experts rather than 73.1\%, but the two variants fail on different tasks: RTE for the full variant and MRPC for the intermediate-only one. Small labelled budgets are unstable across seeds. Limitations include a single set of experts, small validation sets for two tasks, and development data that exceeds the few-shot budget.

\vspace{3mm}
\noindent\textbf{Keywords:} model consolidation; Weight Ensembling; Mixture of Experts; BERT; multi-task learning; few-shot fine-tuning.
